% Part: lambda-calculus
% Chapter: introduction
% Section: minimization

\documentclass[../../../include/open-logic-section]{subfiles}

\begin{document}

\olfileid{lam}{int}{min}
\olsection{Fungsi kang \usetoken{S}{lambda definable} Katutup tumrap Minimisasi}

\begin{lem}
Upamana $f(x,y)$ iku !!{lambda definable}.
Tetepna $g$ kanthi
\[
g(x) \simeq \umin{y}{f(x,y)}.
\]
Mula $g$ uga !!{lambda definable}.
\end{lem}

\begin{proof}
Gagasane kira-kira kaya mangkene. Yen
diwenehi $x$, kita bakal nggunakake
suku lambda titik tetep $Y$ kanggo
netepake fungsi $h_x(n)$ kang nggoleki~$y$
wiwit saka~$n$; banjur $g(x)$ mung
$h_x(0)$. Fungsi~$h_x$ bisa diwujudake
minangka solusi persamaan titik tetep:
\[
h_x(n) \simeq
\begin{cases}
n & \text{yen $f(x,n) = 0$} \\
h_x(n+1) & \text{yen ora.}
\end{cases}
\]

Mangkene rinciane. Amarga $f$ iku
!!{lambda defined} dening sawijining
suku $F$.
Elinga yen kita uga duwe suku
lambda~$D$, saengga
$D(M, N, \bar{0}) \red M$ lan
$D(M, N, \bar{1}) \red N$.
Saiki kanggo sawetara wektu tetepna
$x$; kanggo !!{lambda define} $h_x$
kita kudu nemokake suku~$H$
(kang gumantung marang~$x$)
kang nyukupi
\[
H(\num{n}) \equiv D(\num{n}, H(\fn{Succ}(\num{n})), F(x, \num{n})).
\]
Iki bisa ditindakake nganggo suku
titik tetep~$Y$. Sepisanan, tetepna
$U$ minangka suku
\[
\lambd[h][\lambd[z][D(z,(h(\fn{Succ}(z))),F(x,z))]],
\]
banjur tetepna $H$ minangka suku~$YU$.
Gatekna yen siji-sijine variabel bebas
ing~$H$ yaiku~$x$. Kita bakal nuduhake
yen $H$ nyukupi persamaan ing ndhuwur.

Miturut definisi $Y$, kita duwe
\[
H = YU \equiv U(YU) = U(H).
\]
Mligine, kanggo saben wilangan asli~$n$,
kita duwe
\begin{align*}
H(\num{n}) & \equiv U(H, \num{n}) \\
& \red D(\num{n}, H(\fn{Succ}(\num{n})), F(x, \num{n})),
\end{align*}
kaya kang dikarepake. Gatekna yen
nalika numeral $\num{m}$ disubstitusi
kanggo~$x$ ing larik pungkasan,
ekspresi kasebut direduksi dadi
$\num{n}$ yen $F(\num{m},
\num{n})$ direduksi dadi $\num{0}$,
lan direduksi dadi
$H(\fn{Succ}(\num{n}))$ yen
$F(\num{m}, \num{n})$ direduksi
dadi numeral liyane.

Kanggo ngrampungake bukti, tetepna
$G$ minangka $\lambd[x][H(\num 0)]$.
Banjur $G$ !!{lambda define}s~$g$;
tegese, kanggo saben~$m$,
$G(\num m)$ direduksi dadi
$\overline {g(m)}$ yen $g(m)$
nduweni nilai, lan ora nduweni
wujud normal yen ora.
\end{proof}

\end{document}
